\section{Computational details}\label{app:computation}

This appendix collects the recursions, quadrature rules, simulation algorithms and error
calculations behind \cref{sec:estimation}. The notation is that of
\cref{sec:est-mle,sec:est-em}. Exponential kernels
$g_{kj}(u)=\beta_{kj}e^{-\beta_{kj}u}$ and exponential shock responses
$\psi_k(u)=\eta_ke^{-\eta_ku}$ are assumed unless stated otherwise. The methods are
implemented in the accompanying code (\texttt{code/hawkes}).

\subsection{Recursions for exponential kernels}\label{app:comp-rec}

Fix a pair $(k,j)$. Let $s_1<s_2<\cdots$ be the merged event times of streams $j$ and
$k$ only, $d_m$ the stream of $s_m$, and $\Delta_m=s_m-s_{m-1}$. Define the decay sums
\[
  R_{kj}(t)=\sum_{l:\,t^j_l<t}e^{-\beta_{kj}(t-t^j_l)},\qquad
  R^{(1)}_{kj}(t)=\sum_{l:\,t^j_l<t}(t-t^j_l)\,e^{-\beta_{kj}(t-t^j_l)} .
\]
Write $t-t^j_l=\Delta_m+(s_{m-1}-t^j_l)$ for $t=s_m$, and note that the incident at
$s_{m-1}$ contributes lag $\Delta_m$. This gives
\begin{equation}\label{eq:app-rec}
  R_{kj}(s_m)=e^{-\beta_{kj}\Delta_m}\big\{R_{kj}(s_{m-1})+\ind{d_{m-1}=j}\big\},\qquad
  R^{(1)}_{kj}(s_m)=e^{-\beta_{kj}\Delta_m}R^{(1)}_{kj}(s_{m-1})+\Delta_m\,R_{kj}(s_m),
\end{equation}
with $R_{kj}(s_1)=R^{(1)}_{kj}(s_1)=0$. These two recursions give every quantity that
\cref{sec:estimation} needs at the incidents $t^k_i$.
\begin{itemize}
  \item \emph{Likelihood.} The intensity $\lambda_k(t^k_i)$ is given by
    \eqref{eq:est-lambda} and the compensator by \eqref{eq:est-comp}.
  \item \emph{Compensator at the incidents} (for time rescaling). Since
    $\sum_{t^j_l<t}\{1-e^{-\beta_{kj}(t-t^j_l)}\}=N_j(t-)-R_{kj}(t)$,
    \[
      \Lam_k(t)=\int_0^t\mu_k+\upsilon_k\sum_eS_e\{\Psi_k(t-\tau_e)-\Psi_k(-\tau_e)\}
      +\sum_j\alpha_{kj}\{N_j(t-)-R_{kj}(t)\}.
    \]
  \item \emph{Gradient.} With $\lambda=\lambda_k(t^k_i)$ and sums over the incidents of
    stream $k$,
    \begin{align*}
      \frac{\partial\ell_T}{\partial\alpha_{kj}}&=\sum_i\frac{\beta_{kj}R_{kj}(t^k_i)}{\lambda}-G_{kj},\\
      \frac{\partial\ell_T}{\partial\beta_{kj}}&=\alpha_{kj}\sum_i\frac{R_{kj}(t^k_i)-\beta_{kj}R^{(1)}_{kj}(t^k_i)}{\lambda}
        -\alpha_{kj}\sum_l(T-t^j_l)e^{-\beta_{kj}(T-t^j_l)},\\
      \frac{\partial\ell_T}{\partial\theta_k}&=\sum_i\frac{\mu_k(t^k_i)B(t^k_i)}{\lambda}-\int_0^T\mu_k(t)B(t)\,\dd t,\\
      \frac{\partial\ell_T}{\partial\upsilon_k}&=\sum_i\frac{\sum_eS_e\psi_k(t^k_i-\tau_e)}{\lambda}-\sum_eS_e\Psi^T_k(e).
    \end{align*}
    The derivative in $\eta_k$ uses the shock analogues of $R$ and $R^{(1)}$ below.
  \item \emph{EM statistics.} $N_{kj}$ and $W_{kj}$ of \eqref{eq:est-stats} are the sums
    over $i$ of $\alpha_{kj}\beta_{kj}R_{kj}(t^k_i)/\lambda$ and
    $\alpha_{kj}\beta_{kj}R^{(1)}_{kj}(t^k_i)/\lambda$.
  \item \emph{Shocks.} The shock terms use the same recursions, with the calendar merged
    into the event sequence. The decay sum is
    $R^{\mathrm{sh}}_k(t)=\sum_{\tau_e<t}S_ee^{-\eta_k(t-\tau_e)}$, which jumps by $S_e$
    after $\tau_e$. For short calendars, direct summation over $e$ is simpler.
\end{itemize}

\paragraph{Numerical stability and vectorisation.} The sequential recursion multiplies by
factors $e^{-\beta\Delta_m}\le1$ and is stable. Our implementation vectorises it. Source
events are cut into chunks that span at most $200/\beta_{kj}$ time units. Within a chunk,
$R_{kj}$ is a cumulative sum of $e^{\beta_{kj}(t^j_l-a)}$ relative to an anchor $a$ at
the start of the chunk. A short loop over chunks carries the state forward. The result is
exact up to rounding, with relative error of order $10^{-13}$; there is no truncation
window.

\paragraph{Ties.} All sums use the strict past. If two incidents share a time stamp, the
recursion with $\Delta_m=0$ would count one as the parent of the other. Ties are therefore
broken first, by randomisation within the recorded resolution (\cref{sec:plan-sens}).

\paragraph{Cost.} The recursion for pair $(k,j)$ costs $O(n_j+n_k)$, so all pairs together
cost $O(Mn_{\mathrm{ev}})$. Running it over all $n_{\mathrm{ev}}$ merged events for every pair, as one
state vector updated at every incident, would cost $O(M^2n_{\mathrm{ev}})$.

\paragraph{Parallelism.} If no parameter is shared across rows of the model, the
likelihood separates by target stream (\cref{sec:est-mle}). Each row can then be fitted
on its own processor.

\subsection{Background integrals}\label{app:comp-bg}

The log-linear background needs $\int_0^Te^{B(t)^{\top}\theta}\,\dd t$, together with its
gradient $\int Be^{B^{\top}\theta}$ and Hessian $\int BB^{\top}e^{B^{\top}\theta}$. We
use a composite 8-point Gauss--Legendre rule and reuse the same nodes for all three
integrals. The cells must be chosen so that the integrand is smooth on each of them. For
cubic B-splines alone, the exponent is a polynomial on each knot interval, and our
implementation splits each knot interval into 16 cells. Day-of-week and holiday indicators
make the integrand jump at day boundaries. With such indicators, the cells must also be
split at every day boundary, giving at least one cell per day; our implementation
currently supports smooth bases only. Seasonal (Fourier) bases use uniform cells. The
basis evaluated at the incidents costs $O(n_{\mathrm{ev}}K_b)$ per gradient evaluation.
The M-step for $\theta_k$ (\cref{prop:est-mstep}(a)) takes Newton steps with gradient
$\sum_ip^{k,0}_iB(t^k_i)-\int B\mu_k-\lambda^{\mathrm{pen}}_kP\theta_k$ and negative
Hessian $\int BB^{\top}\mu_k+\lambda^{\mathrm{pen}}_kP$. Step halving ensures that the
M-step objective increases, so the iteration remains a generalised EM and
\cref{prop:est-mstep}(d) applies.

\subsection{Simulation}\label{app:comp-sim}

Synthetic data for the simulation study, and parametric bootstrap replicates, are drawn by
one of two exact algorithms.

\paragraph{Ogata's thinning.} Between incidents and away from calendar times, the
excitation and shock terms of $\lambda_k$ are nonincreasing, because they decay
exponentially. A time-varying background need not be monotone, but it is bounded by
$\bar\mu_k\ge\sup_{[0,T]}\mu_k$. So the sum of $\bar\mu_k$ and the current excitation and
shock terms bounds the future intensity until the next incident or calendar time. \Cref{alg:ogata} exploits this. Candidates come from a Poisson
process at the current bound and are accepted with probability $\lambda/\bar\lambda$.
By the thinning theorem \citep{lewis1979,ogata1981}, the accepted points have intensity
$\lambda$. After a rejection the bound is recomputed at the new time, which is valid
because exponential waiting times are memoryless. At a calendar time the shock term jumps
up, so the algorithm restarts there.

\begin{algorithm}[t]
\caption{Ogata's thinning for the shock-augmented model with exponential kernels.}\label{alg:ogata}
\begin{algorithmic}[1]
\Require $\vartheta$; calendar $\{(\tau_e,S_e)\}$; horizon $T$; bounds
  $\bar\mu_k\ge\sup_{[0,T]}\mu_k$.
\Ensure Events $\mathcal S$ on $[0,T]$ with intensity \eqref{eq:est-lambda}.
\State $t\gets0$; $R_{kj}\gets0$ for all $(k,j)$; $\mathcal S\gets\emptyset$.
\Loop
  \State $h_k\gets\upsilon_k\sum_{\tau_e\le t}S_e\eta_ke^{-\eta_k(t-\tau_e)}$ for all $k$;\quad
    $\bar\lambda\gets\sum_k\{\bar\mu_k+h_k+\sum_j\alpha_{kj}\beta_{kj}R_{kj}\}$.
  \State Draw $E\sim\mathrm{Exp}(\bar\lambda)$; $s\gets t+E$; let $\tau^+$ be the first
    calendar time in $(t,T]$ (or $+\infty$).
  \If{$\tau^+\le s$}
    \State $R_{kj}\gets R_{kj}e^{-\beta_{kj}(\tau^+-t)}$ for all $(k,j)$; $t\gets\tau^+$;
      \textbf{next iteration}.
  \EndIf
  \If{$s>T$} \State \Return $\mathcal S$ \EndIf
  \State $R_{kj}\gets R_{kj}e^{-\beta_{kj}(s-t)}$ for all $(k,j)$; $t\gets s$.
  \State $\lambda_k\gets\mu_k(t)+\upsilon_k\sum_{\tau_e<t}S_e\eta_ke^{-\eta_k(t-\tau_e)}
    +\sum_j\alpha_{kj}\beta_{kj}R_{kj}$ for all $k$; draw $U\sim U(0,1)$.
  \If{$U\bar\lambda\le\sum_k\lambda_k$}
    \State Draw $k$ with probability $\lambda_k/\sum_{k'}\lambda_{k'}$; add $(t,k)$ to
      $\mathcal S$; $R_{k'k}\gets R_{k'k}+1$ for all $k'$.
  \EndIf
\EndLoop
\end{algorithmic}
\end{algorithm}

\paragraph{Exact cluster simulation.} The cluster representation (\cref{sec:th-cluster})
gives a second exact algorithm. It is our default, because it is faster and records the
true branching labels.
\begin{enumerate}[(1)]
  \item \emph{Background immigrants.} Draw a Poisson process with intensity $\mu_k$ on
    $[0,T]$ for each stream. A piecewise-constant $\mu_k$ is sampled directly; any other
    $\mu_k$ by thinning against $\bar\mu_k$.
  \item \emph{Shock immigrants.} For each $(e,k)$, draw $\Pois\{\upsilon_kS_e\Psi^T_k(e)\}$
    points at $\tau_e$ plus delays from $\psi_k$, truncated to the window. The delays are
    drawn by inverting the truncated distribution function.
  \item \emph{Offspring.} Proceed generation by generation. Each point of stream $j$ at
    $s$ has $\Pois(\alpha_{kj})$ children in each stream $k$, placed at $s$ plus
    independent $g_{kj}$ delays. Children after $T$ are discarded; all their descendants
    would also lie after $T$.
\end{enumerate}
When $\spr(\Amat)<1$, the generations form a subcritical multitype Galton--Watson process.
The algorithm then terminates almost surely, after an expected $\Mmat\bm m$ points, where
$\bm m$ is the vector of expected immigrant counts. Each generation is a vectorised
operation, so the cost is $O(Mn_{\mathrm{ev}})$. For a stationary start, either simulate on $[-T_0,T]$
and discard $[-T_0,0)$, or use perfect simulation \citep{moller2005}. The burn-in error
decays with the cascade memory $1/\{\beta(1-n)\}$ (\cref{prop:th-spectrum}). Comparing the
two algorithms on the same parameters is a useful unit test, since they share no code
beyond the model specification.

\subsection{Standard errors}\label{app:comp-se}

Standard errors come from the observed information
$\hat{\bm I}_{\mathrm{obs}}=-\nabla^2\ell_T(\hat\vartheta)$. Our implementation computes it
by central differences of the exact log-likelihood, with relative step $10^{-4}$.
Central differences of the analytic gradient of \cref{app:comp-rec} are more accurate.
They need $2p$ gradient evaluations for $p$ free parameters and have $O(h^2)$ error in each
entry; one-sided differences need only $p$ evaluations but have $O(h)$ error. Two further
details matter.
\begin{itemize}
  \item \emph{Boundary parameters.} Free parameters estimated on or near the boundary
    ($\hat\alpha_{kj}\approx0$, $\hat\upsilon_k\approx0$) are excluded from the inversion and
    flagged. Their sampling distribution is not normal (\cref{prop:lrt}(b)). Intervals for
    them come from the profile likelihood, the parametric bootstrap, or the posterior.
  \item \emph{Alternatives.} The information can also be assembled from complete-data
    quantities via Louis's identity \citep{louis1982}. Two further consistent forms are
    positive semidefinite by construction: the optional (quadratic) variation of the
    score, $\sum_k\sum_i\lambda_k^{-2}\nabla\lambda_k\nabla\lambda_k^{\top}(t^k_i)$, a sum
    over incidents; and its predictable variation,
    $\sum_k\int_0^T\lambda_k^{-1}\nabla\lambda_k\nabla\lambda_k^{\top}\,\dd t$.
\end{itemize}
Derived quantities use the delta method. For the multiplier,
$\partial\Mmat/\partial\alpha_{kj}=\Mmat\bm e_k\bm e_j^{\top}\Mmat$ (\cref{eq:at-grad}).
If $\spr(\Amat)$ is a simple eigenvalue (for example, if $\Amat$ is irreducible) with
right and left Perron vectors $\bm v$ and $\bm w$, then
$\partial\spr(\Amat)/\partial\alpha_{kj}=w_kv_j/(\bm w^{\top}\bm v)$. Near criticality and
near the boundary the delta method is unreliable (\cref{sec:at-uncertainty}), and we use
bootstrap or posterior draws instead.

\subsection{Tests and calibration}\label{app:comp-tests}

\paragraph{Likelihood-ratio tests.} The alternative is fitted first, with the default
multistart. The null fit fixes $\alpha_{kj}=0$ exactly. It is started both from the
alternative optimum, with $\alpha_{kj}$ set to $0$, and from the default multistart, and
the better null optimum is kept. Because the alternative nests the null, a null
log-likelihood above the alternative's signals a local optimum, and the alternative is then
refitted from the null optimum. The statistic is $\max\{0,2(\hat\ell_1-\hat\ell_0)\}$. For
the sup-statistic, each alternative fit fixes $\beta_{kj}$ at a grid value and starts from
the null fit with a small positive $\alpha_{kj}$.

\paragraph{Score process for the Davies problem.} At the null fit $\tilde\vartheta$, the
score in the nuisance parameters vanishes when they are interior. The standardised
efficient score at decay $\beta$ is therefore
\[
  \hat Z(\beta)=\frac{U_\alpha(\beta)}{\{I_{\alpha\alpha}(\beta)-I_{\alpha2}(\beta)I_{22}^{-1}I_{2\alpha}(\beta)\}^{1/2}},\qquad
  U_\alpha(\beta)=\sum_i\frac{\beta R_{kj}(t^k_i;\beta)}{\tilde\lambda_k(t^k_i)}-G_{kj}(\beta).
\]
The information blocks are those of the optional-variation form in
\cref{app:comp-se}, with $\partial\lambda_k/\partial\alpha_{kj}=\beta R_{kj}(\cdot;\beta)$.
Each grid point costs one recursion \eqref{eq:app-rec}. The whole path, and with it the
sup-score statistic and the bound $p_{\mathrm D}$ of \cref{sec:est-lrt}, needs a single
null fit. The same holds in every bootstrap replicate, which makes the score-based
bootstrap cheaper than the likelihood-ratio version by a factor of about the grid size.

\paragraph{Placebo calendars.} Each placebo calendar requires one refit of the selected
model with an extra shock class. The refits are independent and run in parallel. Starting
each refit from the fit to the real calendar, with the placebo loading set to a small
positive value, reduces the number of EM iterations substantially.

\subsection{Complexity and memory}\label{app:comp-cost}

\Cref{tab:app-cost} summarises the costs. Here $n_{\mathrm{ev}}$ is the number of incidents, $M$ the
number of streams, $E$ the number of calendar events, $K_b$ the background basis size,
$Q$ the number of quadrature nodes, $p$ the number of free parameters, $G$ the size of a
decay grid, and $B$ the number of bootstrap or placebo replicates. No step stores an
$n_{\mathrm{ev}}\times n_{\mathrm{ev}}$ array. Direct summation over pairs of incidents, by contrast,
needs $O(n_{\mathrm{ev}}^2)$ time, and $O(n_{\mathrm{ev}}^2)$ memory if the parent probabilities are
stored.

\begin{table}[t]
\centering
\caption{Computational cost of the main steps (exponential kernels unless stated).}\label{tab:app-cost}
\small
\begin{tabularx}{\textwidth}{@{}Xll@{}}
\toprule
Step & Time & Memory\\
\midrule
Log-likelihood or gradient & $O(Mn_{\mathrm{ev}}+{n_{\mathrm{ev}}}E+n_{\mathrm{ev}}K_b+QK_bM)$ & $O(n_{\mathrm{ev}}+M^2+QK_b)$\\
EM iteration (E-step, M-step root solves) & $O(Mn_{\mathrm{ev}}+{n_{\mathrm{ev}}}E+n_{\mathrm{ev}}K_b+r\,Mn_{\mathrm{ev}})$ & $O(n_{\mathrm{ev}}+M^2)$\\
Background Newton step (stream $k$) & $O\{(n_k+Q)K_b^2\}$ & $O\{(n_k+Q)K_b\}$\\
Gibbs sweep (labels, conjugate updates) & $O(Mn_{\mathrm{ev}}+n_{\mathrm{ev}}\log n_{\mathrm{ev}})$ & $O(n_{\mathrm{ev}}+M^2)$\\
Histogram or spline kernels (any shape) & $O(n_{\mathrm{ev}}w)$ per EM iteration & $O(n_{\mathrm{ev}})$\\
Ogata thinning (\cref{alg:ogata}) & $O\{(M^2+E)\,n_{\mathrm{cand}}\}$ & $O(n_{\mathrm{ev}}+M^2)$\\
Cluster simulation & $O(Mn_{\mathrm{ev}})$ expected & $O(n_{\mathrm{ev}})$\\
Observed information by differences & $O(p^2)$ likelihood evaluations & $O(p^2)$\\
Sup-LR test, bootstrap & $O(BG)$ fits & as a fit\\
Sup-score test, bootstrap & $O\{B(\text{null fit}+GMn_{\mathrm{ev}})\}$ & as a fit\\
INAR($p_{\ell}$) by CLS, $N_{\mathrm{bin}}$ bins & $O(p_\ell N_{\mathrm{bin}}M^2+p_\ell^3M^3)$ & $O(p_\ell^2M^2)$\\
\bottomrule
\end{tabularx}
\par\smallskip
\begin{minipage}{\textwidth}\footnotesize
$w$: mean number of incidents within the kernel support. $r$: root-search steps per
M-step. $n_{\mathrm{cand}}$: number of
thinning candidates, $n_{\mathrm{ev}}\int\bar\lambda/\int\lambda$ in expectation. $p_\ell$: number of
INAR lags.
\end{minipage}
\end{table}

\FloatBarrier
